\begin{theorem}[Alling 与 Greenleaf, [1]]\label{thm:2.1.15}
设 $\Sigma$ 是克莱因曲面。则存在 Real 黎曼面 $(\tilde\Sigma,\tau)$ 以及克莱因曲面同态
$p:\tilde\Sigma\to\Sigma$，使得 $p\circ\tau=p$，且诱导映射
$\bar p:\tilde\Sigma/\tau\xrightarrow{\ \simeq\ }\Sigma$ 是克莱因曲面的同构。三元组
$(\tilde\Sigma,\tau,p)$ 称为 $\Sigma$ 的\emph{解析 double}。
\end{theorem}

若 $(\tilde\Sigma',\tau',p')$ 是 $\Sigma$ 的另一个解析 double，则存在解析同构
$f:\tilde\Sigma'\to\tilde\Sigma$ 使得 $f\circ\tau'=\tau\circ f$，即
$(\tilde\Sigma,\tau)$ 与 $(\tilde\Sigma',\tau')$ 作为 Real 黎曼面同构；并且这一构造定义了
克莱因曲面范畴与 Real 黎曼面范畴之间的一个范畴等价。
\begin{annote}
原文此处作「存在解析同构 $f:\Sigma\to\tilde\Sigma$ 使得 $f\circ\tau=\tau'\circ f$，
即 $(\Sigma,\tau)$ 与 $(\Sigma',\tau')$ 同构」，其中 $f$ 的定义域与所用记号都与上下文
不合；按其本意应为 $f:\tilde\Sigma'\to\tilde\Sigma$ 且 $f\circ\tau'=\tau\circ f$。译文
按本意写出，特此说明（参见疑点清单 Q005）。
\end{annote}

最后一句断言的意思是：给定 Real 黎曼面 $(\Sigma,\tau)$，商 $\Sigma/\tau$ 是以
$(\Sigma,\tau)$ 为解析 double 的克莱因曲面；并且两个克莱因曲面 $\Sigma_1,\Sigma_2$ 之间的
同态与它们解析 double 之间的 Real 黎曼面同态一一对应。注意克莱因曲面的解析 double 是
定义 2.1.7 意义下的 double（我们已回忆过，本讲义中的 Real 黎曼面边界为空）。若克莱因
曲面 $\Sigma$ 事实上是（无边）黎曼面，则其解析 double 为
$\tilde\Sigma=\Sigma\sqcup\bar\Sigma$，其中 $\bar\Sigma$ 表示把 $\Sigma$ 上的全纯函数
换为反全纯函数所得的曲面（特别地，底集上的恒等映射 $\Sigma\to\bar\Sigma$ 是反全纯的，
因而定义了 $\tilde\Sigma$ 上的一个实结构）。[1] 中还介绍了其他的 double 概念（定向
double 与 Schottky double）。

要勾勒给定克莱因曲面 $\Sigma$ 的解析 double 的存在性证明并不太难。若
$(\phi_i:U_i\xrightarrow{\simeq}V_i)_{i\in I}$ 是 $\Sigma$ 的双解析图册，则在 $\C$ 中
作开集 $W_i=V_i\cup\bar V_i$，并考虑拓扑空间 $\Omega:=\bigsqcup_{i\in I}W_i$。为简单
起见，假定所有 $V_i$ 都是 $\C^+$ 中的圆盘或半圆盘。若 $\phi_i\circ\phi_j^{-1}$ 解析，
则把 $\phi_j(U_i\cap U_j)\subset V_j$ 粘到 $\phi_i(U_i\cap U_j)\subset V_i$，并把
$\phi_j(U_i\cap U_j)\subset\bar V_j$ 粘到 $\phi_i(U_i\cap U_j)\subset\bar V_i$；若
$\phi_i\circ\phi_j^{-1}$ 反解析，则把 $\phi_j(U_i\cap U_j)\subset V_j$ 粘到
$\phi_i(U_i\cap U_j)\subset\bar V_i$，并把 $\phi_j(U_i\cap U_j)\subset\bar V_j$ 粘到
$\phi_i(U_i\cap U_j)\subset V_i$。这就在 $\Omega$ 上定义了一个等价关系 $\sim$，令
$\tilde\Sigma:=\Omega/\sim$ 并赋予商拓扑。用 $[W_i]$ 记 $W_i\subset\Omega$ 在典范投影
下的像。存在典范同胚 $z_i:[W_i]\to W_i$，并且注意 $\tilde\Sigma$ 的边界为空。由 $\Omega$
上粘合的定义，映射 $z_i\circ z_j^{-1}$ 总是解析的：例如若 $\phi_i\circ\phi_j^{-1}$
反解析，则 $z_i\circ z_j^{-1}$ 是映射
\[
v\longmapsto
\begin{cases}
\phi_i\circ\phi_j^{-1}(v), & \text{若 } v\in V_j,\\[0.3em]
\overline{\phi_i\circ\phi_j^{-1}(\bar v)}, & \text{若 } v\in\bar V_j,
\end{cases}
\]
故由 Schwarz 反射原理它是解析的（真正用到反射原理的只是 $V_i\cap V_j$ 与 $\C^+$ 中的
$\R$ 相交的情形）。$\tilde\Sigma$ 上的实结构 $\tau$ 在局部坐标卡 $([W_i],z_i)$ 中就是
复共轭；而若在局部解析坐标卡中读出，投影 $p:\tilde\Sigma\to\Sigma$ 把
$x=a+\sqrt{-1}\,b\in W_i$ 送到 $a+\sqrt{-1}\,|b|$，因此它是克莱因曲面同态。定理 2.1.15
证明的其余部分可参看 [1]。

\begin{remark}\label{rem:2.1.16}
在本讲义余下的部分，我们只关心连通的 Real 黎曼面 $(\Sigma,\tau)$，即克莱因曲面
$\Sigma/\tau$ 的底拓扑曲面要么不可定向，要么边界非空（两个性质可以同时成立）。我们将
不再处理双解析结构，因而常把对子 $(\Sigma,\tau)$ 本身称为克莱因曲面。
\end{remark}

\subsection*{习题}

\begin{exercise}\label{ex:2.1.17}
回忆 $\Aut(\CP^1)\simeq\mathrm{PGL}(2;\C)$，它以分式线性变换作用。通过分析在该群中
「是二阶元」意味着什么，证明 $\CP^1$ 上仅有的可能实结构是 $z\mapsto\bar z$ 与
$z\mapsto-1/\bar z$。
\end{exercise}

\begin{exercise}\label{ex:2.1.18}
设 $a,b,c$ 是三个互不相同的复数，且集合 $\{a,b,c\}\subset\C$ 在 $\C$ 的复共轭下不变。
考虑由方程 $P(x,y)=0$ 定义的实椭圆曲线，其中 $P(x,y)=y^2-(x-a)(x-b)(x-c)$。
\begin{enumerate}[label=(\roman*), leftmargin=2.2em, itemsep=0.15em]
\item 证明：若 $a,b,c$ 是实数，则由 $P$ 定义的实平面曲线有两个连通分支（在 $\R^2$ 的
      通常 Hausdorff 拓扑下）。
\item 证明：若例如 $a$ 不是实数，则由 $P$ 定义的实平面曲线只有一个连通分支。
\item 在上述两种情形下，由 $P$ 定义的实射影曲线的拓扑型各是什么？
\end{enumerate}
\end{exercise}

\begin{exercise}\label{ex:2.1.19}
设 $P\in\C[x,y]$ 是二元复系数非常数多项式。证明
$X_P(\C)=\{(u,v)\in\C^2\mid P(u,v)=0\}$ 是非紧黎曼面。
\end{exercise}

\begin{exercise}\label{ex:2.1.20}
证明 (2.1) 中定义的层 $(p_*\calO_\Sigma)^\tau$ 是 $\Sigma/\tau$ 上的 $\R$-代数层，
并且经 (2.2) 它是 $\Sigma/\tau$ 上光滑函数层的子层。
\end{exercise}

\begin{exercise}\label{ex:2.1.21}
\begin{enumerate}[label=(\roman*), leftmargin=2.2em, itemsep=0.15em]
\item 证明：若 $(P,\sigma)$ 是可定向曲面连同其一个反转定向对合，则 $P^\sigma$ 是至多
      $g+1$ 个圆周之并。（提示：只需证明 $P^\sigma$ 是一维流形，并重述定理 2.1.5 的
      证明；要点在于不用复解析结构与全纯映射的规范形来证明命题 2.1.3。）
\item 证明：相反地，可定向曲面的保定向对合的不动点集是零维的。
\end{enumerate}
\end{exercise}

\begin{exercise}\label{ex:2.1.22}
对曲面 $S_1,S_2$ 的连通和，证明经典关系
$\chi(S_1\#S_2)=\chi(S_1)+\chi(S_2)-2$。并重新检查定理 2.1.5 的证明。
\end{exercise}

\begin{exercise}\label{ex:2.1.23}
画出亏格 $g$ 的定向曲面 $P$ 的一个反转定向对合 $\sigma$ 的图形，使 $\sigma$ 恰有
$g+1$ 个实点圆周。
\end{exercise}

\begin{exercise}\label{ex:2.1.24}
设 $\Sigma$ 是亏格 $g$ 的紧连通黎曼面。
\begin{enumerate}[label=(\roman*), leftmargin=2.2em, itemsep=0.15em]
\item 证明 $\Sigma$ 上实结构的拓扑型有 $\bigl[(3g+4)/2\bigr]$ 个。
\item 对较小的 $g$ 计算这个数，并在每种情形下从拓扑上识别曲面 $\Sigma/\tau$。
      （提示：用 Weichold 定理 2.1.8，并按 $g$ 的奇偶分情形讨论。）
\end{enumerate}
\end{exercise}

\begin{exercise}\label{ex:2.1.25}
\begin{enumerate}[label=(\roman*), leftmargin=2.2em, itemsep=0.15em]
\item 证明：定义在闭上半平面 $\C^+$ 的开子集上、满足 Schwarz 反射原理（定理 2.1.11）
      假设 (i)--(iii) 的函数构成 $\R$-代数层。
\item 证明：如此定义的局部环层空间 $(A,\mathcal{D}_A)$ 同构于 $(U/\sigma,(p_*\calO_U)^\sigma)$，
      其中 $U=A\cup\bar A$，$\sigma$ 是复共轭，$(p_*\calO_U)^\sigma$ 如习题 2.1.20 所定义。
\end{enumerate}
\end{exercise}

\begin{exercise}\label{ex:2.1.26}
设 $A$ 是闭上半平面 $\C^+$ 的开子集。证明双解析函数 $f:A\to\C^+$ 是克莱因曲面同态。
\end{exercise}

\begin{exercise}\label{ex:2.1.27}
满足 $k=g+1$ 的实曲线 $(\Sigma,\tau)$ 称为\emph{极大曲线}（或 \emph{M-曲线}）。证明
$k=g+1$ 当且仅当 $\Sigma^\tau$ 的 Betti 数之和（不要忘记对 $\Sigma^\tau$ 的 $g+1$ 个
连通分支求和）等于 $\Sigma$ 的 Betti 数之和。取模 $2$ 系数，这就给出把 Harnack 不等式
推广到高维（光滑射影）实代数簇的办法：实部的模 $2$ Betti 总数小于或等于复部的模 $2$
Betti 总数；这称为 Milnor--Thom 不等式（[12, 19]）。
\end{exercise}

\section{实代数曲线的基本群}
\label{sec:02.2}

\subsection{黎曼面基本群简忆}
\label{sec:02.2.1}

连通带基点拓扑空间 $(X,x)$ 的拓扑基本群 $\pi_1(X;x)$ 通常取为连续映射
$\gamma:[0,1]\to X$（满足 $\gamma(0)=\gamma(1)=x$，也称为 $x$ 处的环路）的同伦类构成的
群。下文我们常直接把环路 $\gamma$ 的同伦类也记作 $\gamma$（这应不至于引起混淆）。这个群
的一个关键性质是：它左作用于任一拓扑覆盖映射 $p:Y\to X$ 在 $x$ 处的纤维。事实上，这是
下述引理的推论（本节自始至终假定 $X$ 连通，且 $X,Y$ 是 Hausdorff 空间）。

\begin{lemma}[道路与同伦的提升]\label{lem:2.2.1}
设 $\gamma:[0,1]\to X$ 是从 $x$ 出发的道路（即满足 $\gamma(0)=x$ 的连续映射
$\gamma:[0,1]\to X$），$p:Y\to X$ 是拓扑覆盖映射。则
\begin{enumerate}[label=(\roman*), leftmargin=2.2em, itemsep=0.15em]
\item 给定 $y\in p^{-1}(\{x\})$，存在唯一的道路 $\tilde\gamma:[0,1]\to Y$，使得
      $p\circ\tilde\gamma=\gamma$ 且 $\tilde\gamma(0)=y$；
\item 若 $\gamma'$ 同伦于 $\gamma$，则 $\tilde\gamma'$ 同伦于 $\tilde\gamma$；特别地，
      $\tilde\gamma'(1)=\tilde\gamma(1)$。
\end{enumerate}
\end{lemma}

引理 2.2.1 的证明例如可参看 [5]。条件 $p\circ\tilde\gamma=\gamma$ 说明：给定
$y\in p^{-1}(\{x\})$ 与 $x$ 处的环路 $\gamma$，元素
\begin{equation}\label{eq:2.5}
\gamma\curvearrowright y:=\tilde\gamma(1)
\end{equation}
属于 $p^{-1}(\{x\})$，且 $\gamma\curvearrowright y$ 只依赖于 $\gamma$ 的同伦类。本讲义
采用的道路复合约定是
\begin{equation}\label{eq:2.6}
(\gamma*\gamma')(t)=
\begin{cases}
\gamma'(2t), & \text{若 } 0\le t\le\tfrac12,\\[0.3em]
\gamma(2t-1), & \text{若 } \tfrac12\le t\le1,
\end{cases}
\end{equation}
（即先走 $\gamma'$，再走 $\gamma$），因此 (2.5) 确实定义了 $\pi_1(X;x)$ 对任一以 $X$ 为
底的拓扑覆盖在 $x$ 处纤维的左作用。

\begin{definition}[单值作用]\label{def:2.2.2}
给定以 $X$ 为底的拓扑覆盖 $p:Y\to X$、点 $x\in X$ 与点 $y\in p^{-1}(\{x\})$，由提升
$x$ 处环路所定义（即 $\tilde\gamma$ 是 $Y$ 中满足 $p\circ\tilde\gamma=\gamma$、
$\tilde\gamma(0)=y\in Y$ 的唯一道路）的左作用
$\gamma\curvearrowright y:=\tilde\gamma(1)$，称为\emph{单值作用}（monodromy action）。
\end{definition}

单值作用在下述意义下是函子式的：给定 $X$ 的两个覆盖 $p:Y\to X$ 与 $q:Z\to X$，
以及它们之间的同态
\[
\begin{tikzcd}[column sep=small, row sep=small]
Y \arrow[rr, "\phi"] \arrow[dr, "p"'] & & Z \arrow[dl, "q"]\\
 & X &
\end{tikzcd}
\]
则诱导映射 $\phi_x:p^{-1}(\{x\})\to q^{-1}(\{x\})$ 是 $\pi_1(X;x)$-等变的。事实上，若
$\tilde\gamma$ 满足 $p\circ\tilde\gamma=\gamma$ 且 $\tilde\gamma(0)=y\in p^{-1}(\{x\})$，
则 $\phi\circ\tilde\gamma$ 满足
\[
q\circ(\phi\circ\tilde\gamma)=(q\circ\phi)\circ\tilde\gamma
=p\circ\tilde\gamma=\gamma,
\]
且 $(\phi\circ\tilde\gamma)(0)=\phi(y)$，故
\begin{equation}\label{eq:2.7}
\gamma\curvearrowright\phi(y)=(\phi\circ\tilde\gamma)(1)
=\phi\bigl(\tilde\gamma(1)\bigr)=\phi(\gamma\curvearrowright y).
\end{equation}
这实际上把 $\pi_1(X;x)$ 变成了「把 $X$ 的覆盖送到其在 $x$ 处纤维」这一函子的自同构群。
然而，对于连通、Hausdorff 且局部单连通的拓扑空间 $X$，万有覆盖的存在大大简化了事情。

\begin{definition}[万有覆盖]\label{def:2.2.3}
设 $X$ 是连通拓扑空间，$p:\tilde X\to X$ 是拓扑覆盖映射。称 $p$ 为 $X$ 的\emph{万有
覆盖}，若 $\tilde X$ 连通，且对 $x\in X$、$\tilde x\in\tilde X$，下述泛性质成立：
对任何连通覆盖 $q:Y\to X$ 及任意选定的 $y\in p^{-1}(\{x\})$，存在唯一的拓扑覆盖映射
$r:\tilde X\to Y$，使得 $q\circ r=p$ 且 $q(\tilde x)=y$：
\[
\begin{tikzcd}[column sep=large, row sep=large]
(\tilde X,\tilde x) \arrow[dr, "p"'] \arrow[rr, "\exists!\,r"] & & (Y,y) \arrow[dl, "q"]\\
 & (X,x) &
\end{tikzcd}
\]
\end{definition}

重要的是要注意，上述泛性质是带基点空间的同态 $p:(\tilde X,\tilde x)\to(X,x)$ 所具有的
性质。此后我们总假定万有覆盖映射是带基点空间的同态，即使记号没有明确反映已选定点
$\tilde x\in p^{-1}(\{x\})$。我们要指出，把 $p:\tilde X\to X$ 按通常方式\footnote{回忆
构造的一部分是给 $\tilde X$ 赋予使 $p$ 连续的拓扑，在该拓扑下 $\tilde X$ 实际上是连通
的——特别地，它只能覆盖 $X$ 的其他连通覆盖。}构造为从 $x$ 出发的道路
$\eta:[0,1]\to X$ 的同伦类之集、$p$ 为把 $\eta$ 送到 $\eta(1)$ 的映射，这依赖于
$x\in X$ 的选取，但它随后定义了一个特殊的点 $\tilde x\in p^{-1}(\{x\})$，即 $x$ 处
常道路的同伦类。此外，在这个模型中 $\pi_1(X;x)$ 自然地作用在 $\tilde X$ 上，这是一个
右作用：
\begin{equation}\label{eq:2.8}
\forall\eta\in\tilde X,\ \forall\gamma\in\pi_1(X;x),\qquad
\eta\curvearrowleft\gamma:=\eta*\gamma
\end{equation}
（回忆：先走 $\gamma$ 再走 $\eta$，故 $\eta*\gamma$ 确实是 $X$ 中从 $x$ 出发的道路）。
显然 (2.8) 定义了 $\pi_1(X;x)$ 在 $\tilde X$ 上的右作用，且是覆盖映射
$p:\tilde X\to X$ 的自同构作用，因为
$p(\eta*\gamma)=(\eta*\gamma)(1)=\eta(1)=p(\eta)$（回忆覆盖 $p:\tilde X\to X$ 的自同构群
为 $\Aut(\tilde X/X):=\{f:\tilde X\to\tilde X\mid p\circ f=p\}$，群结构由映射复合给出）。
注意这一构造定义了从 $\pi_1(X;x)$ 到 $\Aut(\tilde X/X)$ 的映射。在命题 2.2.4 中，我们
不用 $x$ 处环路来构造同一个映射。

由于泛性质，带基点万有覆盖 $p:(\tilde X,\tilde x)\to(X,x)$（若存在）在同构意义下唯一，
而且同构唯一。有一个定理说：若 $p:\tilde X\to X$ 是 Hausdorff 空间之间的拓扑覆盖映射，
且 $\tilde X$ 连通且单连通，则 $(\tilde X,\tilde x)$ 是 $(X,x)$ 的万有覆盖。此外，若 $X$
连通、Hausdorff 且（半）局部单连通（例如连通拓扑流形），则上面勾勒的构造所得的
$\tilde X$ 连通、Hausdorff 且单连通，故原空间 $X$ 特别地具有万有覆盖（例如可参看 [5]）。
现在我们放下 $\tilde X$ 的这个特殊模型，直接假定万有覆盖
$p:(\tilde X,\tilde x)\to(X,x)$ 存在，并已知 $\tilde X$ 单连通。这使我们能把
$\pi_1(X;x)$ 描述为 $\Aut(\tilde X)^{\mathrm{op}}$（$\tilde X$ 的自同胚群，群结构由
$f\cdot g=g\circ f$ 给出）的子群，比进入函子自同构群的形式体系更简单。回忆：若 $G$ 是
以 $\cdot$ 为复合律的群，则 $G^{\mathrm{op}}$ 是在 $G$ 的底集上由规定 $g$ 与 $h$ 在
$G^{\mathrm{op}}$ 中的乘积 $g\cdot^{\mathrm{op}}h$ 等于 $h\cdot g$ 而得的群结构；它称为
$G$ 的\emph{反群}。

\begin{proposition}\label{prop:2.2.4}
由万有覆盖的泛性质，给定 $\gamma\in\pi_1(X;x)$，存在唯一的 $f:\tilde X\to\tilde X$，
它在 $X$ 上覆盖 $\Id_X$ 且满足 $f(\tilde x)=\gamma\curvearrowright\tilde x$（即 (2.5)
中定义的单值作用）。这实际上定义了群同构
\begin{equation}\label{eq:2.9}
\Psi:\ \pi_1(X;x)\xrightarrow{\ \simeq\ }\Aut(\tilde X/X)^{\mathrm{op}},
\qquad \gamma\longmapsto f.
\end{equation}
\end{proposition}

\begin{proof}
该映射是良定义的。我们先验证它是群同态，然后给出逆映射。记 $\Psi(\gamma)=f$，
$\Psi(\gamma')=g$。按定义，$\Psi(\gamma*\gamma')$ 是唯一的 $h:\tilde X\to\tilde X$，它在
$X$ 上覆盖 $\Id_X$ 且满足 $h(\tilde x)=(\gamma*\gamma')\curvearrowright\tilde x$。但
\[
(\gamma*\gamma')\curvearrowright\tilde x
=\gamma\curvearrowright(\gamma'\curvearrowright\tilde x)
=\gamma\curvearrowright g(\tilde x),
\]
而 $g$ 对单值作用是 $\pi_1(X;x)$-等变的（见 (2.7)），故
\[
h(\tilde x)=\gamma\curvearrowright g(\tilde x)
=g(\gamma\curvearrowright\tilde x)=g\bigl(f(\tilde x)\bigr)
=(g\circ f)(\tilde x)=(f\cdot g)(\tilde x).
\]
由这种映射的唯一性，$\Psi(\gamma*\gamma')=f\cdot g=\Psi(\gamma)\cdot\Psi(\gamma')$，
故 $\Psi:\pi_1(X;x)\to\Aut(\tilde X/X)^{\mathrm{op}}$ 是群同态。其次，考虑把
$f\in\Aut(\tilde X/X)^{\mathrm{op}}$ 送到 $\gamma:=p\circ\zeta$ 的同伦类的映射，其中
$\zeta$ 是 $\tilde x$ 与 $f(\tilde x)\in p^{-1}(\{x\})$ 之间的任一道路：由于 $\tilde X$
单连通，$\zeta$ 的同伦类唯一确定，从而 $\gamma$ 的同伦类是良定义的 $\pi_1(X;x)$ 中
元素，不依赖于 $\zeta$ 的选取。此外，$\zeta$ 是从 $\tilde x$ 出发的 $\gamma$ 的提升，
故 $\gamma\curvearrowright\tilde x=\zeta(1)=f(\tilde x)$，这直接说明此构造给出 $\Psi$ 的
逆映射。
\end{proof}

若接受单连通万有覆盖的存在性，命题 2.2.4 的结果可用作 $\pi_1(X;x)$ 的定义。必须考虑
$\Aut(\tilde X/X)^{\mathrm{op}}$ 而非 $\Aut(\tilde X/X)$，根源在于 (2.6) 中为道路复合
所选定的约定，而该约定实际上是由局部系统理论决定的（见 [18, Remark 2.6.3]）。特别地，
$\Aut(\tilde X/X)^{\mathrm{op}}$ 右作用于 $\tilde X$。下面我们说明它左作用于任一覆盖
$q:Y\to X$ 的纤维 $q^{-1}(\{x\})$。这只需用到 $p:(\tilde X,\tilde x)\to(X,x)$ 的泛性质，
并给出单值作用 (2.5) 的一个不使用 $x$ 处环路的描述。
